%%%%%%%%%%%%%%%%%%%%%%%%%%%%%%%%%%%%%%%%%%%%%%%%%%%%%%%%%%%%%%%%%%%%%%
% Lớp 11 | Đạo hàm | Các quy tắc đạo hàm
% Dạng 1D7?2-6: Giới hạn hàm số lượng giác, hàm số mũ, lôgarít
% Số câu: 17  (sinh tự động bởi tach_cau.py; từ nay thầy sửa trực tiếp file này)
%%%%%%%%%%%%%%%%%%%%%%%%%%%%%%%%%%%%%%%%%%%%%%%%%%%%%%%%%%%%%%%%%%%%%%

%%%%%%%%%%%%%%%% TRẢ LỜI NGẮN %%%%%%%%%%%%%%%%

% nguồn: Ok_TraLoiNgan_Lop11_Lan1
\begin{ex}%[1D7V2-6]%[ma:89effb5e]
	Cho hàm số $f(x)=\sin 2x+\sqrt{3}\cos 2x+12\sin \left(x+\dfrac{\pi}{6}\right)$. Họ nghiệm của phương trình $f'(x)+4=0$ có dạng  $x=\dfrac{\pi}{m}+k\pi$ với $m$ nguyên dương, $k\in \mathbb{Z}$. Giá trị của $m$ bằng
	 \shortans{3}
	\loigiai{ 
		\begin{itemize}
			\item $f'(x)=2\cos 2x-2\sqrt{3}\sin 2x+12\cos\left(x+\dfrac{\pi}{6}\right).$ 
			 \item $f'(x)+4=0\Leftrightarrow \dfrac{1}{2}\cos 2x-\dfrac{\sqrt{3}}{2}\sin 2x+3\cos \left(x+\dfrac{\pi}{6}\right)+1=0$\\
			$\Leftrightarrow \cos\left(2x+\dfrac{\pi}{3}\right)+3\cos\left(x+\dfrac{\pi}{6}\right)+1=0\Leftrightarrow 2\cos^2\left(x+\dfrac{\pi}{6}\right)+3\cos\left(x+\dfrac{\pi}{6}\right)=0$\\ 
			$\Leftrightarrow \cos\left(x+\dfrac{\pi}{6}\right)=0$ (vì  $\cos\left(x+\dfrac{\pi}{6}\right)\in [-1;1]$).\\
			$\Leftrightarrow x+\dfrac{\pi}{6}=\dfrac{\pi}{2}+k\pi\Leftrightarrow x=\dfrac{\pi}{3}+k\pi, k\in\mathbb{Z}$.\\ 
			 Vậy $m=3$.
		\end{itemize}
	 }
\end{ex}


%%%%%%%%%%%%%%%% TỰ LUẬN %%%%%%%%%%%%%%%%

% nguồn: Ok_TuanLuan_Toan11_Lan1
\begin{bt}%[1D7H2-6]%[ma:4a8f5e6e]
	Tính đạo hàm  của các hàm số sau
	\begin{listEX}[2]
		\item $y=2x^3-\dfrac{5}{x}+4 \sqrt{x}-1$. 
		\item $y=x \cdot \log_{2}x$. 
	\end{listEX}
	\loigiai{ 
		\begin{enumerate}
			\item  $y'=6x^2+\dfrac{5}{x^2}+\dfrac{4}{2 \sqrt{x}}$.
			\item $y'= x' \cdot \log_{2}x+x \cdot (\log_{2}x)' =  \log_{2}x + x \cdot \dfrac{1}{x \ln2} = \log_{2}x + \dfrac{1}{\ln 2}$.
		\end{enumerate}
	 }
\end{bt}

% nguồn: Ok_TuanLuan_Toan11_Lan2
\begin{ex}%[1D7H2-6]%[ma:bd20d6f8]
	Tính đạo  hàm của các hàm số sau
	\begin{listEX}[2]
		\item  $y=2x^3-\dfrac{5}{x}+4 \sqrt{x}-1$.
		\item  $y=x \cdot \log_{2}x$.
	\end{listEX}
	 \loigiai{
		\begin{enumerate}
			\item  $y'=6x^2+\dfrac{5}{x^2}+\dfrac{4}{2 \sqrt{x}}$.
			\item $y'= x' \cdot \log_{2}x+x  \cdot (\log_{2}x)' = \log_{2}x + x \cdot \dfrac{1}{x \ln2} = \log_{2}x + \dfrac{1}{\ln 2}$.
		\end{enumerate}
	 }
\end{ex}

% nguồn: Ok_TuanLuan_Toan11_Lan2
\begin{ex}%[1D7V2-6]%[ma:d9c4b3b3]
	Tìm giới hạn $L=\lim\limits_{x \to 0} \dfrac{\ln  (\sqrt[3]{3 x+1}+1)-\ln (\sqrt{x+1}+1)}{x}$.
	\loigiai{\allowdisplaybreaks 
		\begin{eqnarray*}
			& L&=\lim\limits_{x \to 0} \dfrac{\ln \left(\sqrt[3]{3  x+1}+1\right)-\ln \left(\sqrt{x+1}+1\right)}{x}\\
			& &=\lim\limits_{x \to 0} \dfrac{\ln \left(\sqrt[3]{3 x+1}+1\right)-\ln 2-\left[\ln  \left(\sqrt{x+1}+1\right)-\ln 2\right]}{x}\\
			& &=\lim\limits_{x \to  0}\left[\dfrac{\ln \left(1+\dfrac{\sqrt[3]{3 x+1}-1}{2}\right)}{x}-\dfrac{\ln \left(\dfrac{\sqrt{x+1}-1}{2}+1\right)}{x}\right].
		\end{eqnarray*}
		 Ta có\allowdisplaybreaks
		\begin{eqnarray*}
			&&\lim\limits_{x \to 0} \dfrac{\ln \left(1+\dfrac{\sqrt[3]{3 x+1}-1}{2}\right)}{x}\\ 
			&=&\lim\limits_{x \to 0} \dfrac{\ln \left(1+\dfrac{1}{2}\cdot \dfrac{3 x}{\sqrt[3]{(3 x+1)^2}+\sqrt[3]{3 x+1}+1}\right)}{\dfrac{2}{3}\cdot \left(\sqrt[3]{(3  x+1)^2}+\sqrt[3]{3 x+1}+1\right) \cdot \dfrac{1}{2}\cdot \dfrac{3 x}{\left(\sqrt[3]{(3 x+1)^2}+\sqrt[3]{3 x+1}+1\right)}}\\
			&=&\dfrac{1}{2}. 
		\end{eqnarray*}
		Và $\lim\limits_{x \to 0} \dfrac{\ln \left(\dfrac{\sqrt{x+1}-1}{2}+1\right)}{x}=\lim\limits_{x \to 0} \dfrac{\ln \left(\dfrac{1}{2} \cdot  \dfrac{x}{\sqrt{x+1}+1}+1\right)}{2\left(\sqrt{x+1}+1\right)\cdot \dfrac{1}{2} \cdot \dfrac{x}{\sqrt{x+1}+1}}=\dfrac{1}{4}$.\\
		Vậy $L=\dfrac{1}{2}-\dfrac{1}{4}=\dfrac{1}{4}$. 
	 }
\end{ex}

% nguồn: Ok_TuanLuan_Toan11_Lan2
\begin{ex}%[1D7V2-6]%[ma:ceeff90d]
	Tìm giới hạn $L=\lim\limits_{x \to  0} \dfrac{\mathrm{e}^{-2 x^2}-\sqrt[3]{1+x^2}}{\ln \left(1+x^2\right)}$.
	\loigiai{ Ta có 
	\begin{eqnarray*}
		&  L&=\lim\limits_{x \to 0} \dfrac{\mathrm{e}^{-2 x^2}-\sqrt[3]{1+x^2}}{\ln \left(1+x^2\right)}\\
		& &=\lim\limits_{x \to 0}\left(\dfrac{\mathrm{e}^{-2 x^2}-1}{x^2}-\dfrac{\sqrt[3]{1+x^2}-1}{x^2}\right): \dfrac{\ln  \left(1+x^2\right)}{x^2} \\
		& &=\lim\limits_{x  \to 0}\left(-2\cdot \dfrac{\mathrm{e}^{-2 x^2}-1}{-2 x^2}-\dfrac{1}{\sqrt[3]{\left(1+x^2\right)^2}+\sqrt[3]{1+x^2}+1}\right): \dfrac{\ln \left(1+x^2\right)}{x^2}=-\dfrac{7}{3}.
	\end{eqnarray*}
	 }
\end{ex}

% nguồn: Ok_TuanLuan_Toan11_Lan2
\begin{ex}%[1D7V2-6]%[ma:632d2048]
	Tìm giới hạn $A=\lim\limits_{x \to 0} \dfrac{\mathrm{e}^{x^2}-\cos x}{x^2}$. 
	 \loigiai{Ta có
	\begin{eqnarray*}
		& A&=\lim\limits_{x \to 0}  \dfrac{\mathrm{e}^{x^2}-\cos x}{x^2}\\
		& &=\lim\limits_{x \to 0}\left(\dfrac{\mathrm{e}^{x^2}-1}{x^2}+\dfrac{1-\cos x}{x^2}\right)=\lim\limits_{x \to 0}\left(\dfrac{\mathrm{e}^{x^2}-1}{x^2}+\dfrac{2 \sin ^2 \dfrac{x}{2}}{4 \cdot\left(\tfrac{x}{2}\right)^2}\right). 
	\end{eqnarray*}
	 Suy ra $A=1+\dfrac{1}{2}=\dfrac{3}{2}$.\\
	Vậy $A=\dfrac{3}{2}$. 
	 }
\end{ex}

% nguồn: Ok_TuanLuan_Toan11_Lan2
\begin{ex}%[1D7V2-6]%[ma:173fa369]
	Tìm giới hạn  $A=\lim\limits_{x \to 0} \dfrac{\mathrm{e}^{\sqrt{2 x+1}-1}-\mathrm{e}^{\sqrt[3]{1-3 x}-1}}{x}$.
	 \loigiai{
	 $$
	A=\lim\limits_{x \to 0} \dfrac{\mathrm{e}^{\sqrt{2 x+1}-1}-\mathrm{e}^{\sqrt[3]{1-3 x}-1}}{x}=\lim\limits_{x \to 0}\left(\dfrac{\sqrt{2 x+1}-1}{x} \cdot  \dfrac{\mathrm{e}^{\sqrt{2 x+1}-1}-1}{\sqrt{2 x+1}-1}-\dfrac{\sqrt[3]{1-3 x}-1}{x} \cdot \dfrac{\mathrm{e}^{\sqrt[3]{1-3 x}-1}-1}{\sqrt[3]{1-3 x}-1}\right).
	$$ 
	Vì $\lim\limits_{x \to 0}  \dfrac{\sqrt{2 x+1}-1}{x}=\lim\limits_{x \to 0} \dfrac{2 x}{x(\sqrt{2 x+1}+1)}=\lim\limits_{x \to 0} \dfrac{2}{\sqrt{2 x+1}+1}=1$ nên
	 $$\lim\limits_{x \to 0} \left(\dfrac{\sqrt{2 x+1}-1}{x} \cdot \dfrac{\mathrm{e}^{\sqrt{2 x+1}-1}-1}{\sqrt{2 x+1}-1}\right)=1.$$
	Ta có $\lim\limits_{x \to 0} \dfrac{\sqrt[3]{1-3 x}-1}{x}=\lim\limits_{x \to 0} \dfrac{-3  x}{x\left[\sqrt[3]{(1-3 x)^2}+\sqrt[3]{1-3 x}+1\right]}=\lim\limits_{x \to 0} \dfrac{-3}{\left[\sqrt[3]{(1-3 x)^2}+\sqrt[3]{1-3 x}+1\right]}=-1$ nên
	$$\lim\limits_{x \to 0} \left( \dfrac{\sqrt[3]{1-3 x}-1}{x} \cdot \dfrac{\mathrm{e}^{\sqrt[3]{1-3  x}-1}-1}{\sqrt[3]{1-3 x}-1}\right)=-1.$$
	 Vậy $A=2$.
	 }
\end{ex}

% nguồn: Ok_TuanLuan_Toan11_Lan2
\begin{ex}%[1D7V2-6]%[ma:3b8f8864]
	Tìm giới hạn $A=\lim\limits_{x  \to+\infty}\left(\dfrac{1+x}{x-1}\right)^x$.
	\loigiai{Với $A=\lim\limits_{x \to+\infty}\left(\dfrac{1+x}{x-1}\right)^x$, đặt $t=x-1$, khi $x \to+\infty$ thì  $t \to+\infty$. Do đó
	$$ 
	A=\lim\limits_{t \to+\infty}\left(\dfrac{t+2}{t}\right)^{t+1}=\lim\limits_{t \to+\infty} \dfrac{t+2}{t}\left(\dfrac{t+2}{t}\right)^t=\lim\limits_{t \to+\infty}  \dfrac{t+2}{t}\cdot \left[\left(1+\dfrac{2}{t}\right)^{\tfrac{t}{2}}\right]^2=\mathrm{e}^2.
	 $$
	 Vậy $A=\mathrm{e}^2$.
	} 
\end{ex}

% nguồn: Ok_TuanLuan_Toan11_Lan2
\begin{ex}%[1D7V2-6]%[ma:0987a4d4]
	Cho $ a \in \mathbb{N}^* $, $ a \le 9 $. Tìm giới hạn $A=\lim  \dfrac{a+\overline{aa}+\cdots+\underbrace{\overline{aa\ldots a}}_{ n \text{ chữ số } a}}{10^n}$.
	\loigiai{ Ta  có
	\begin{eqnarray*}
		& a+\overline{aa}+\cdots+\underbrace{\overline{aa\ldots a}}_{ n \text{ chữ số } a}&=a\cdot (1+11+111+\cdots+\underbrace{111  \ldots 1}_{n \text { chữ số } 1})\\
		& &=\dfrac{a}{9}\cdot (9+99+999+\ldots+\underbrace{999 \ldots  9}^{n \text { chữ số }9})\\
		 & & =\dfrac{a}{9}\cdot\left(10-1+10^2-1+10^3-1+\cdots+10^n-1\right)\\
		 & &=\dfrac{a}{9}\cdot\left[\dfrac{10\left(10^n-1\right)}{9}-n\right].
	\end{eqnarray*}
	Ta có $\lim \dfrac{a+\overline{aa}+\cdots+\underbrace{\overline{aa\ldots a}}_{ n \text{ chữ số } a}}{10^n}=\lim \dfrac{a}{9\cdot 10^n}\left[\dfrac{10\left(10^n-1\right)}{9}-n\right]=\lim \dfrac{a}{9}\left[\dfrac{10\left(10^n-1\right)}{9\cdot  10^n}-\dfrac{n}{10^n}\right]$.\\
	 Mặt khác $\lim \dfrac{10\left(10^n-1\right)}{9\cdot 10^n}=\dfrac{10}{9}$.\\
	 Và $n \leq \mathrm{C}_n^0+\mathrm{C}_n^1+\cdots+\mathrm{C}_n^n=2^n \Rightarrow 0 \leq \dfrac{n}{10^n} \leq\left(\dfrac{2}{10}\right)^n=\dfrac{1}{5^n}$, mà $\lim \dfrac{1}{5^n}=0 \Rightarrow \lim \dfrac{n}{10^n}=0$.\\
	 Vậy $A=\dfrac{10 a}{81}$.
	 }
\end{ex}

% nguồn: Ok_TuanLuan_Toan11_Lan2
\begin{ex}%[1D7V2-6]%[ma:9acb40f7]
	Tìm giới hạn  $A=\lim\limits_{x \to 0} \dfrac{\mathrm{e}^{a x}-\mathrm{e}^{b x}}{x}$ $\left(a,b\in\mathbb{R}^*\right)$.
	 \loigiai{
	Ta có $A=\lim\limits _{x \to 0} \dfrac{\mathrm{e}^{a x}-\mathrm{e}^{b x}}{x}=\lim\limits _{x \to  0}\left(a\cdot \dfrac{\mathrm{e}^{a x}-1}{a x}-b\cdot \dfrac{\mathrm{e}^{b x}-1}{b x}\right)=a-b$.\\
	Vậy $ A=a-b  $
	} 
\end{ex}

% nguồn: Ok_TuanLuan_Toan11_Lan2
\begin{ex}%[1D7V2-6]%[ma:71ba4e83]
	 Tìm giới hạn $L=\lim\limits_{x \to 0} \dfrac{\ln (\sin x+\cos x)}{x}$.
	\loigiai{Ta  có
	\begin{eqnarray*}
		&L&=\lim\limits_{x \to 0} \dfrac{\ln (\sin x+\cos  x)}{x}\\
		&  &=\lim\limits_{x \to 0} \dfrac{2 \ln (\sin x+\cos x)}{2 x}\\
		& &=\lim\limits_{x \to 0} \dfrac{\ln (\sin x+\cos  x)^2}{2 x}\\
		& &=\lim\limits_{x \to 0} \dfrac{\ln (1+\sin 2 x)}{2 x}\\ 
		& &=\lim\limits_{x \to 0}\left[\dfrac{\ln (1+\sin 2 x)}{\sin 2 x} \cdot \dfrac{\sin 2 x}{2  x}\right]=1.
	\end{eqnarray*}
	} 
\end{ex}

% nguồn: Ok_TuanLuan_Toan11_Lan2
\begin{ex}%[1D7V2-6]%[ma:3aa9cd24]
	Tìm giới hạn $A=\lim\limits_{x \to 0} \dfrac{(1+x)^\alpha-1}{x}$. 
	\loigiai{ 
	Ta có $$ 
	A=\lim\limits_{x \to 0} \dfrac{(1+x)^\alpha-1}{x}=\lim\limits_{x \to 0} \alpha\cdot \dfrac{\ln (1+x)}{x} \cdot \dfrac{\mathrm{e}^{\alpha  \ln (1+x)}-1}{\alpha \ln (1+x)}=\alpha.$$ Vậy $A=\alpha$.
	 }
\end{ex}

% nguồn: Ok_TuanLuan_Toan11_Lan2
\begin{ex}%[1D7V2-6]%[ma:c20bbb69]
	Tìm  giới hạn $A=\lim\limits_{x \to 0} \dfrac{\mathrm{e}^{-3 x^2}-\sqrt[3]{1+x^2}}{\ln \left(1+x^2\right)}$.
	\loigiai{Ta có 
	$$ 
	A=\lim\limits_{x \to 0} \dfrac{\mathrm{e}^{-3 x^2}-\sqrt[3]{1+x^2}}{\ln \left(1+x^2\right)}=\lim\limits_{x \to  0}\left[(-3)\cdot \dfrac{\mathrm{e}^{-3 x^2}-1}{\left(-3 x^2\right)} \cdot \dfrac{x^2}{\ln \left(1+x^2\right)}+\dfrac{1-\sqrt[3]{1+x^2}}{\ln \left(1+x^2\right)}\right].
	$$ 
	 Ta có
	\begin{itemize}
		\item $\lim\limits_{x \to 0}\left((-3) \cdot \dfrac{\mathrm{e}^{-3  x^2}-1}{\left(-3 x^2\right)} \cdot \dfrac{x^2}{\ln \left(1+x^2\right)}\right)=-3$.
		\item $\lim\limits_{x \to 0}  \dfrac{1-\sqrt[3]{1+x^2}}{\ln \left(1+x^2\right)}=\lim\limits_{x \to 0}\left[\dfrac{-1}{1+\sqrt[3]{1+x^2}+\left(\sqrt[3]{1+x^2}\right)^2} \cdot \dfrac{x^2}{\ln \left(1+x^2\right)}\right]=-\dfrac{1}{3}$.
	\end{itemize}
	Vậy  $A=-3-\dfrac{1}{3}=-\dfrac{10}{3}$.
	 }
\end{ex}

% nguồn: Ok_TuanLuan_Toan11_Lan2
\begin{ex}%[1D7V2-6]%[ma:b96d4f49]
	Tìm giới hạn $L=\lim\limits_{x \to 0}  \dfrac{6^x-3^x}{\ln (1+6 x)-\ln (1+3 x)}$.
	\loigiai{Ta có 
	\begin{eqnarray*}
		&L=\lim\limits_{x \to 0} \dfrac{6^x-3^x}{\ln (1+6 x)-\ln (1+3 x)}&=\lim\limits_{x \to 0}\left(\dfrac{6^x-1}{x}-\dfrac{3^x-1}{x}\right):\left(\dfrac{\ln  (1+6 x)}{x}-\dfrac{\ln (1+3 x)}{x}\right)\\
		 & &=\dfrac{\ln 6-\ln 3}{6-3}=\dfrac{1}{3} \ln 2.
	\end{eqnarray*}
	 }
\end{ex}

% nguồn: Ok_TuanLuan_Toan11_Lan2
\begin{ex}%[1D7V2-6]%[ma:97dcd120]
	Tìm giới hạn $A=\lim\limits_{x \to 0} \dfrac{1-\mathrm{e}^x}{\sqrt{x+1}-1}$. 
	 \loigiai{ Ta có
	 $$
	A=\lim\limits_{x \to 0} \dfrac{1-\mathrm{e}^x}{\sqrt{x+1}-1}=\lim\limits_{x  \to 0} \left(\dfrac{-x}{\sqrt{x+1}-1} \cdot \dfrac{\mathrm{e}^x-1}{x}\right).
	$$ 
	Vì $\lim\limits_{x \to 0} \dfrac{-x}{\sqrt{x+1}-1}=-\lim\limits_{x \to  0}(\sqrt{x+1}+1)=-2$ nên $\lim\limits_{x \to 0} \left( \dfrac{-x}{\sqrt{x+1}-1} \cdot \dfrac{\mathrm{e}^x-1}{x}\right)=-2$.\\
	Vậy  $A=-2$.
	 }
\end{ex}

% nguồn: Ok_TuanLuan_Toan11_Lan2
\begin{ex}%[1D7V2-6]%[ma:c6c9620a]
	Tìm giới hạn $A=\lim\limits_{x \to 0}  \dfrac{\sqrt{\mathrm{e}^x}-1}{\sin 2 x}$.
	 \loigiai{ Ta có
	$$ 
	A=\lim\limits_{x \to 0} \dfrac{\sqrt{\mathrm{e}^x}-1}{\sin 2 x}=\lim\limits_{x \to 0} \left(  \dfrac{\mathrm{e}^x-1}{x} \cdot \dfrac{1}{2\left(\sqrt{\mathrm{e}^x}+1\right)} \cdot \dfrac{2 x}{\sin 2 x}\right)=\dfrac{1}{4}.$$
	 Vậy $A=\dfrac{1}{4}$.
	 }
\end{ex}

% nguồn: Ok_TuanLuan_Toan11_Lan2
\begin{ex}%[1D7V2-6]%[ma:cde91e1e]
	Tìm giới hạn $L=\lim\limits_{x \to 0} \dfrac{\ln  \left(1+3 x^2\right)}{1-\cos 2 x}$.
	 \loigiai{
	Ta có $L=\lim\limits_{x \to 0} \dfrac{\ln \left(1+3 x^2\right)}{1-\cos 2 x}=\lim\limits_{x \to 0}  \dfrac{\ln \left(1+3 x^2\right)}{2 \sin ^2 x}=\dfrac{1}{2} \lim\limits_{x \to 0}\left[\dfrac{3 \ln \left(1+3 x^2\right)}{3 x^2}:\left(\dfrac{\sin x}{x}\right)^2\right]=\dfrac{3}{2}$.
	} 
\end{ex}
